\section{Matching the primal and dual temperatures}
\label{sec:matching}

Proposition~\ref{prop:two-trajectories} supplies a small trajectory for
the primal couplings at $\beta_{\rm p}$ and another for the dual
couplings at $\beta_{\rm d}$. We now prove that these physical
temperatures coincide. The comparison takes place in the same bond
environment on a torus. Exact duality relates the four spin partition
functions, while FK monotonicity separates two distinct temperatures.
The argument extends the torus identification in
\cite[Section~8, Proposition ``Identification of the tuned parameter'']{R};
it uses the two trajectories in place of symmetry of the disorder law.

\subsection{FK differentiation at fixed environment}

For a finite graph $G=(V,E)$ and a partition $\pi$ of its boundary
vertices, let
\begin{equation}
 \phi_{G,\beta,\omega}^{\pi}(\eta)
 =\frac{1}{\mathcal Z_{G,\beta,\omega}^{\pi}}
       2^{k_\pi(\eta)}\prod_{e\in E}w_e(\beta)^{\eta_e},
 \qquad \eta\in\{0,1\}^{E},
 \label{eq:fk-law}
\end{equation}
where $w_e(\beta)=e^{2\beta J_e}-1$ and $k_\pi$ counts open clusters
after the boundary identifications. The superscript $\mathrm w$
denotes full wiring. All connection events below use ordinary edges;
boundary identifications are not steps in a path. At fixed environment,
these measures are monotonic measures: every edge's conditional open
probability increases when more of the other edges are open. They
satisfy positive association, increase with the edge odds and the
boundary wiring, and have the domain Markov property
\cite{FK,FKG}. Edwards--Sokal coloring couples them to the Ising law
\cite{ES}. The dual measure, denoted by an additional superscript $*$,
has odds $w^*_{e^*}(\beta)=2/w_e(\beta)$ as in
\eqref{eq:dual-coupling}. The parameter $\beta$ in this notation is
always the physical primal inverse temperature.

Two finite-volume formulas will be used both here and in
Section~\ref{sec:criticality}. For an increasing event $A$, set
$q(\beta)=\phi_{G,\beta,\omega}^{\pi}(A)$. Differentiating the finite
sum in \eqref{eq:fk-law} gives
\begin{equation}
 q'(\beta)=\sum_{e\in E}a_e(\beta)
       \Cov_{\phi_{G,\beta,\omega}^{\pi}}(\1_A,\eta_e),
 \qquad
 a_e(\beta)=\frac{\dd}{\dd\beta}\log w_e(\beta)
           =\frac{2J_e}{1-e^{-2\beta J_e}}.
 \label{eq:weighted-derivative}
\end{equation}
On every compact interval in $(0,\infty)$ these coefficients have
uniform positive lower bounds, because
$1-\varepsilon\le J_e\le1+\varepsilon$. For the dual law the same
identity holds with coefficient $-a_e(\beta)$ at the crossing edge.
All the covariances are nonnegative.

If a decision tree determines $A$ and queries edge $e$ with probability
$\delta_e$, the decision-tree inequality for monotonic measures gives
\begin{equation}
 \phi(A)(1-\phi(A))
 \le C_{\rm DT}\sum_{e\in E}\delta_e\Cov_\phi(\1_A,\eta_e).
 \label{eq:fk-osss}
\end{equation}
Here $C_{\rm DT}$ is an absolute constant; an independent randomized
choice of tree is allowed. We use the monotonic-measure extension
\cite[Theorem~1.1]{DCRT} of the OSSS inequality~\cite{OSSS}.
Both \eqref{eq:weighted-derivative} and \eqref{eq:fk-osss} are applied
at fixed environment, before any environmental expectation.

\subsection{Four sectors on a torus}

Torus duality mixes partition functions with different spin twists,
so we retain all four sectors. Normalizing their vector removes a
common bulk factor introduced by the scale transformations. The
resulting comparison will make the primal FK winding probabilities
at the two tuned temperatures asymptotically equal. The derivative bound will
then force those temperatures to agree.

Let $\mathbb T_N=(\Z/N\Z)^2$, $N\ge3$. A flat spin twist with
holonomy $\alpha=(\alpha_1,\alpha_2)\in\{0,1\}^2$ is an assignment
of edge signs $\tau_e^\alpha$ whose product around every elementary
face is $1$ and whose products around the two coordinate cycles are
$(-1)^{\alpha_1}$ and $(-1)^{\alpha_2}$. Define
\[
 Z_\alpha(K)=\sum_{\sigma\in\{-1,1\}^{\mathbb T_N}}
   \exp\left(\sum_{e=\{x,y\}}K_e\tau_e^\alpha\sigma_x\sigma_y\right).
\]
A change of the twist by a vertex gauge transformation is absorbed
by changing the spin variables. Thus $Z_\alpha$ depends only on the
holonomy. Order the four entries as $00,10,01,11$, and write
$Z(K)=(Z_\alpha(K))_\alpha$ and
$\widehat Z(K)=Z(K)/\|Z(K)\|_2$.

\begin{lemma}[Inhomogeneous torus duality]
\label{lem:torus-duality}
For positive couplings $K$ on $\mathbb T_N$, let
$K^*_{e^*}=g(K_e)$ and identify the dual square torus with a translate
of the primal torus. Then
\begin{equation}
 Z(K)=a(K)\mathsf H Z(K^*),\qquad
 a(K)=\prod_e\sqrt{\sinh(2K_e)},\qquad
 \mathsf H=\frac12
 \begin{pmatrix}
 1&1&1&1\\
 1&1&-1&-1\\
 1&-1&1&-1\\
 1&-1&-1&1
 \end{pmatrix}.
 \label{eq:torus-duality}
\end{equation}
The matrix $\mathsf H$ is orthogonal and fixes $Z(K_0)$. Moreover,
under the untwisted torus FK law,
\begin{equation}
 \frac{Z_\alpha(K)}{Z_{00}(K)}
 =\phi_{\mathbb T_N,K}
 \bigl(\alpha\cdot[\gamma]=0
          \text{ for every open cycle }\gamma\bigr),
 \label{eq:torus-ratio}
\end{equation}
where $[\gamma]\in(\Z/2\Z)^2$ is the mod-two winding class.
\end{lemma}

\begin{proof}
This is the finite-volume Kramers--Wannier identity~\cite{KW}; the
four-sector formulation is also given in~\cite{BS} and proved in
\cite[Section~8, Lemma ``The four torus partition functions'']{R}.
We include the expansions to specify the homology and normalization.
For $b\in(\Z/2\Z)^2$, put
\[
 U_b(K)=\sum_{\substack{F\subset E(\mathbb T_N)\ \mathrm{even}\\{}[F]=b}}
                        \prod_{e\in F}\tanh K_e.
\]
Here even means that every vertex has even degree in $F$. The
high-temperature expansion is
\[
 Z_\alpha(K)=2^{N^2}\prod_e\cosh K_e
                   \sum_b(-1)^{\alpha\cdot b}U_b(K).
\]
Since $e^{-2K^*_{e^*}}=\tanh K_e$, the low-temperature expansion of
the dual spin system with holonomy $\gamma$ gives
\[
 Z_\gamma(K^*)=2e^{\sum_eK^*_{e^*}}
                       U_{(\gamma_2,\gamma_1)}(K).
\]
The exchange of coordinates expresses the intersection pairing of
primal and dual cycles; the factor $2$ counts the two spin assignments
with a prescribed domain-wall configuration. Substitution yields the
sign $(-1)^{\alpha_1\gamma_2+\alpha_2\gamma_1}$ in
\eqref{eq:torus-duality}. Its prefactor follows from
$|E(\mathbb T_N)|=2N^2$ and
$e^{-K^*_{e^*}}\cosh K_e=\sqrt{\sinh(2K_e)/2}$.
Direct multiplication gives $\mathsf H^2=I$. At constant coupling
$K_0$, duality leaves the couplings unchanged and $a(K_0)=1$,
so $\mathsf HZ(K_0)=Z(K_0)$.

For the ratio, use
\[
 e^{K_e\tau_e^\alpha\sigma_x\sigma_y}
 =e^{-K_e}\left(1+(e^{2K_e}-1)
             \1_{\{\sigma_x=\tau_e^\alpha\sigma_y\}}\right).
\]
On an open subgraph the constraints are consistent exactly when
every open cycle has zero evaluation under $\alpha$. If consistent,
they admit $2^{k(F)}$ spin assignments. Summation over the subgraphs
and division by the untwisted partition function prove
\eqref{eq:torus-ratio}.
\end{proof}

\begin{proposition}[Matching of the two trajectories]
\label{prop:matching}
The physical temperatures in Proposition~\ref{prop:two-trajectories}
coincide:
\begin{equation}
             \beta_*:=\beta_{\rm p}=\beta_{\rm d}.
 \label{eq:matched-temperature}
\end{equation}
Thus both the primal and the exact dual coupling arrays at $\beta_*$
have uniformly small scale trajectories with $X_k\to0$.
\end{proposition}

\begin{proof}
Choose a fixed sufficiently large integer $M$ and set $N_k=Ms_k$.
The torus is large enough that every complete update stencil used
through level $k$, including field supports, integration boundaries,
and environment-dependence neighborhoods, fits in an embedded square
of diameter less than $N_k/4$. These connected square carriers give
contractible charts with connected nonempty overlaps. Put iid bonds
on this torus and use the same
bond variables at all temperatures being compared.

First consider either of the two trajectories, with couplings
$\widetilde K$. Run its local scale map in all four holonomies,
using the same algorithmic variables, denoted by $\zeta$. In a
contractible chart the flat twist can be removed by a local spin
trivialization. On connected overlaps two such trivializations differ
by a constant spin flip. The potentials are even, and the reference
auxiliary variables are unchanged by spin flip. The local change of
spin coordinates preserves the supremum norms and the flag decisions
made using the shared thresholds, and identifies the twisted pure
conditional kernels with their plane counterparts.
Centering uses the full-plane reference law in these local coordinates,
so its scalar is independent of the sector; it does not use a twisted
torus expectation. The sector-dependent global correlations remain in
the reference partition functions $Z_\alpha(K_0)$. The bare centering has
the same property: writing $c_{\rm en}=\E_0(\sigma_x\sigma_y)$,
the centered twisted bond is
$\tau_e^\alpha\sigma_x\sigma_y-c_{\rm en}$, and the scalar removed
from the initial logarithmic perturbation is
$c_{\rm en}\sum_e(\widetilde K_e-K_0)$ in every holonomy.
Consequently every local integration identity holds with the same
centering factor in all four sectors. This is the all-sector use of
the exact local map, as in the torus argument of~\cite[Section~8]{R}.

After $k$ steps there is therefore a common positive factor
$c_k(\omega,\zeta)$ such that, for every $\alpha$,
\begin{equation}
 e^{-S_k}Z_\alpha(K_0)
 \le c_k(\omega,\zeta)^{-1}Z_\alpha(\widetilde K)
 \le e^{S_k}Z_\alpha(K_0).
 \label{eq:sector-sandwich}
\end{equation}
One may take $S_k$ to be the sum, over the four sectors, of the
supremum norms of all remaining potentials. Integrations in the
reference auxiliary variables are with probability measures and do
not enlarge this bound. There are $M^2$ cells at the final scale,
and embedded stencils have their plane laws. Hence
\[
             \E S_k\le C M^2X_k\longrightarrow0.
\]
Normalizing \eqref{eq:sector-sandwich} in Euclidean norm gives
\begin{equation}
 \bigl\|\widehat Z(\widetilde K)-\widehat Z(K_0)\bigr\|_2
          \longrightarrow0
       \quad\text{in environmental probability}.
 \label{eq:tuned-sector-convergence}
\end{equation}
The partition functions themselves do not involve $\zeta$, so this
conclusion remains true for the original bond environment alone.

For the primal trajectory, \eqref{eq:tuned-sector-convergence}
applies directly at $\beta_{\rm p}$. For the dual trajectory it
applies to $K^*(\beta_{\rm d})$. By
Lemma~\ref{lem:torus-duality},
$\widehat Z(K)=\mathsf H\widehat Z(K^*)$ and the target is fixed
by $\mathsf H$. Thus the \emph{primal} normalized partition vector
approaches the pure vector at both $\beta_{\rm p}$ and
$\beta_{\rm d}$. This uses exact duality in the given environment;
it makes no assertion that the primal and dual environments have
the same distribution.

Take $\alpha=10$, and let $A_\alpha$ be the increasing event that
there is an open cycle with odd evaluation under $\alpha$. Write
$q_{N,\omega}(\beta)=\phi_{\mathbb T_N,\beta,\omega}(A_\alpha)$
and let $q_N^0$ be its pure critical probability. The pure torus
crossing bounds give a constant $c>0$ such that
\begin{equation}
                    c\le q_N^0\le1-c
 \label{eq:pure-torus-nondegeneracy}
\end{equation}
for all sufficiently large $N$. Indeed, a fixed chain of overlapping
FK rectangle crossings constructs a horizontal open winding cycle.
A vertical closed-dual winding cycle excludes every open cycle of
odd horizontal winding. Each construction has a uniform positive
probability by pure FK rectangle bounds, imposed successively under
arbitrary induced wirings; see \cite[Section~8]{R} and
\cite[Section~2, ``Pure buffered RSW consequences'']{S}.

Equation~\eqref{eq:torus-ratio} identifies
$1-q_{N,\omega}=Z_{10}/Z_{00}$. All four partition-function ratios
belong to $[0,1]$, so the untwisted coordinate of every normalized
vector is at least $1/2$. Taking ratios in
\eqref{eq:tuned-sector-convergence} therefore gives, jointly,
\begin{equation}
 q_{N_k,\omega}(\beta_{\rm p})-q_{N_k}^0\longrightarrow0,
 \qquad
 q_{N_k,\omega}(\beta_{\rm d})-q_{N_k}^0\longrightarrow0
 \label{eq:matching-endpoints}
\end{equation}
in environmental probability.

Suppose that $\beta_-:=\min\{\beta_{\rm p},\beta_{\rm d}\}$ is
strictly less than
$\beta_+:=\max\{\beta_{\rm p},\beta_{\rm d}\}$. Querying every
edge in \eqref{eq:fk-osss} and using
\eqref{eq:weighted-derivative} yields, uniformly over this interval,
\[
 q'_{N_k,\omega}(\beta)
       \ge c_1q_{N_k,\omega}(\beta)(1-q_{N_k,\omega}(\beta)).
\]
With probability tending to one, both endpoint values lie in
$[c/2,1-c/2]$ by \eqref{eq:matching-endpoints}. Monotonicity keeps
the whole intervening curve in this interval. Integration then gives
\[
 q_{N_k,\omega}(\beta_+)-q_{N_k,\omega}(\beta_-)
       \ge c_2(\beta_+-\beta_-)>0
\]
on events of probability tending to one. The left side tends to zero
in probability by \eqref{eq:matching-endpoints}, a contradiction.
\end{proof}
